\section{Discussion}
\label{sec:discussion}

\paragraph{Principal findings.} Three findings shape the design. First, check-ins carry most of the information: in simulation, three KSS ratings in the first 16~hours cut the error of the latent sleepiness estimate by about 31\% relative to the rhythm model alone. Second, passive signals add only a little, and how they are added matters more than whether they are added. Fused naively with check-ins, they gave no gain and made the twin overconfident; anchored by per-channel bias states, they were never worse than check-ins alone, kept the 90\% intervals close to nominal, and gave the lowest Brier score for the one-hour forecast of KSS $\geq$ 7. Third, learning each person's slopes from a few check-ins did not help, either in simulation or in DROZY, where even a shrunk intercept shift from one check-in did worse than the population mapping. The public data also show why a single population mapping should not be fully trusted for an individual: in ds004902, objective slowing and subjective sleepiness changed together only weakly across people.

\paragraph{What is new.} BlinkTwin is, to our knowledge, the first open tool that treats validated questionnaires as calibration observations of a personal state-space model that also contains a biomathematical prior and live passive sensing (Table~\ref{tab:comparison}). Individualised fatigue models \citep{rajaraman2008individualized,reifman2019twobalert} adapt a biomathematical model to repeated PVT performance; camera detectors measure the eyes with population thresholds; questionnaires capture the person's own state but only now and then. BlinkTwin joins these in one Kalman filter that keeps each passive channel's personal bias separate from the person's state, gives uncertainty, forecasts the time to KSS~7, answers what-if questions, explains its estimate exactly, and raises real-time alarms, all in a browser with video kept on the device. The same code runs the website, the tests and the simulation, and its blink detector was validated against EOG.

\paragraph{Relation to individualised fatigue models.} The two approaches are complementary. Individualised two-process models predict performance; BlinkTwin predicts the person's own rating of sleepiness and workload. A later version could individualise the prior's parameters as well as the offset and the channel biases.

\paragraph{Limitations.}
\begin{itemize}[leftmargin=*,itemsep=2pt]
  \item The sensing was validated on a different setup from the one the website uses. DROZY used a 10-minute PVT and a near-infrared Kinect camera with 68-point landmarks and the eye aspect ratio; the website uses the 3-minute PVT-B and MediaPipe blendshapes from an RGB webcam. Lighting, glasses, camera angle and screen glare will affect the webcam; the per-person normalisation and the channel bias reduce, but do not remove, these effects. Blendshape eye closure is not clinical PERCLOS.
  \item The blink meter undercounted EOG blinks by 5.4 per minute on average, with wide limits of agreement; blink rate is therefore a relative, not absolute, measure.
  \item Browser reaction times include display and input latency that varies between devices.
  \item The Three-Process Model constants are population values, they still need to be checked against the original sources, and the implementation omits the breakpoint in sleep recovery.
  \item The simulation is only as good as its assumptions. Its person-level distributions came from 14 people in DROZY, the twin's time-on-task and prior parameters are the same in the simulator and the twin, and the workload model was not evaluated.
  \item DROZY is small (14 people, 36 tests), its random-slope models did not converge, and its calibration analysis rests on 22 predictions. In ds004902, KSS and the reaction test were collected on different participants, so person-level coupling was tested with the SSS in 25 people.
  \item The alarm thresholds are defaults that have not been validated.
  \item BlinkTwin is not a medical device. It should not be used to decide whether someone is fit to drive or to operate machinery.
\end{itemize}

\paragraph{Ethics and privacy.} Webcam monitoring of people at work raises clear concerns about surveillance. BlinkTwin is designed for self-monitoring: frames never leave the page and are never stored, data stay in the person's browser, and the person can export or delete them at any time. Any deployment by an employer would need consent, transparency about what is measured, and a guarantee that the data are not used against the person.

\paragraph{Future work.} New signals enter the model as further observations of the same state, each with its own bias. Natural next steps are heart-rate variability from a wearable, eye tracking and wearable EEG, and contexts beyond desk work, such as simulated driving and walking under cognitive load, where cognitive state and motor performance can be measured together. Whether a twin anchored in one context transfers to another is an open question that this framework makes easy to study.

\section{Conclusion}
\label{sec:conclusion}

BlinkTwin is an open, on-device digital twin of sleepiness and cognitive load that runs in a web browser with no installation. Its distinctive features are: validated questionnaires used as calibration observations rather than labels; a bias state for each passive channel, which in simulation turned naive fusion that gave no gain and lost calibration into anchored fusion that was never worse than check-ins alone and gave the best one-hour forecasts; one filter that joins a biomathematical model with live sensing to give uncertainty, a time-to-KSS-7 forecast and what-if scenarios; real-time microsleep and head-drop alarms with video never stored or sent; an exact additive explanation of each estimate; and one tested code base whose blink detector agrees with EOG. Check-ins carried most of the information, and calibration from one or two check-ins did not beat the population mapping in DROZY. Whether personal calibration helps real people is the central question of the pre-specified validation study.
